\documentclass[11pt]{article}

\usepackage[margin=1in]{geometry}
\usepackage{amsmath,amssymb}

\title{A Portability Note for the Local Phase Package Beyond Zeta}
\author{Working note in \texttt{grh/}}
\date{2026-04-23}

\begin{document}

\maketitle

\begin{abstract}
This note records a conservative portability template extracted from the current
RH draft. It does not claim a proof of GRH, ERH, or a Ramanujan-\(\tau\)
analogue. The main conclusion is narrower: the draft appears to contain a
source-light local package based on a real phase, jet normalization,
whitening, an odd transverse scalar template, and an \(N\)-point odd-moment
projector, while the actual zeta-side realization, calibration amplitude,
microscopic denominator control, and late contradiction remain source-specific.
\end{abstract}

\section{Verdict}

The current RH draft appears to split into two layers.

\begin{enumerate}
\item An abstract odd-germ / odd-projector package that is genuinely
source-light.
\item A separate realization and calibration package that is still
source-specific.
\end{enumerate}

The strongest honest theorem boundary is therefore two-layered:

\begin{enumerate}
\item Once a real odd holomorphic transverse scalar with the required
microscopic boundary control exists, the odd expansion and the universal
\(N\)-point odd-moment projector follow.
\item Calibration enters only as a separate conditional interface, after one has
a real sign-compatible amplitude, a realized value-channel derivative, and a
nonzero toy pairing in the same channel.
\end{enumerate}

The safe portability statement is therefore conditional even locally:

\medskip
\noindent
If a completed \(L\)-function supplies the same real critical-line channel,
corrected odd holomorphic transverse scalar, microscopic denominator control,
and preserved whitening scales, then the same odd-projector machinery should
apply. If it also supplies a real sign-compatible amplitude and a nonzero toy
pairing, then the same calibration interface should apply.
\medskip

This note does not claim that those hypotheses have been proved for primitive
Dirichlet \(L\)-functions, for complex characters, or for the Ramanujan
\(\tau\) \(L\)-function.

\section{Portable Local Layer}

The strongest candidate portable layer is:

\begin{enumerate}
\item the phase kernel built from a real phase \(\Phi\);
\item jet normalization and the finite-\(s\) local block package;
\item whitening of same-point and mixed blocks;
\item the odd transverse scalar template;
\item the discrete \(N\)-point odd-moment projector.
\end{enumerate}

In the current draft, these are the parts that look most abstraction-friendly.
The follow-up attack sharpened the stopping point: the unconditional local
theorem should stop at the odd-germ/projector layer.

\section{Conditional Calibration Interface}

The value-channel derivative, the small-\(x\) pairing against the toy anomaly,
and the comparison
\[
u^2\asymp \frac{x}{B}S(m)
\]
should be treated as a second conditional layer, not as part of the
unconditional source-light theorem. Keeping this separation avoids hiding the
family-specific need for a real sign-compatible amplitude and a nonvanishing
pairing denominator.

\section{Nonportable Source Layer}

The first explicitly zeta-side layer begins where the draft introduces the
split
\[
q(t)=B_\zeta(t)+S(t).
\]
From that point on, the argument uses source-specific data:

\begin{enumerate}
\item the real calibration amplitude \(S(m)\) or its replacement;
\item curvature control of the form \(|S''(m)|\ll L(m)S(m)^2\);
\item microscopic denominator comparability for omitted zeros;
\item zeta-side boundary control for the corrected odd scalar;
\item the late RH-specific provenance / endpoint / propagation package.
\end{enumerate}

The research cycle found no support for claiming that this package already
exists outside zeta.

At the same time, a later focused attack sharpened the role of the scalar
\(S(m)\): from the manuscript formulas alone, it is naturally a positive
strip-edge zero kernel,
\[
S(m)=\sum_\rho \Re\left(\frac{1}{1+2im-\rho}-\frac{1}{2im-\rho}\right),
\]
which strongly suggests a completed-object scattering interpretation. This is
best read as a candidate-object refinement, not as a proved formula-level
theorem of the manuscript.

\section{Dirichlet Cases}

For primitive Dirichlet \(L\)-functions, the most plausible first route is to
work on a completed-object package that produces a positive scalar analogue of
\(S(m)\), not directly on raw complex values \(L(\tfrac12+it,\chi)\).

A later focused attack sharpened this further at the phase level: the quotient
\[
\phi_F(s)=\frac{\Lambda_F(2s-1)}{\Lambda_F(2s)}
\]
is naturally unimodular on the critical line under the standard completed-
\(L\) functional equation and conjugation package. So the single completed
quotient is now best viewed as a serious \emph{phase-channel candidate}. That
upgrade is still only a front-end improvement. It does not by itself produce
the positive scalar analogue of \(S(m)\) needed later.

Two cases should be kept separate.

\begin{enumerate}
\item Real or otherwise self-dual cases, where the local phase package may be
closer to the zeta situation.
\item Genuinely non-self-dual complex characters, where a rotated single-channel
realization may still be insufficient and a paired or matrix-valued package may
be needed to recover a positive scalar analogue of \(S(m)\).
\end{enumerate}

The research team and verifier agreed that the non-self-dual case is the least
supported by the present draft.

\section{Ramanujan \texorpdfstring{$\tau$}{tau} Test Case}

The \(\tau\) \(L\)-function looks like a cleaner test object than complex
Dirichlet characters at the self-duality level. The conservative reading is:

\begin{enumerate}
\item degree \(2\) does not obviously break the local matrix algebra;
\item degree \(2\) does change the archimedean background and analytic-conductor
bookkeeping;
\item the present draft does not show that the same microscopic denominator,
boundary, and whitening-scale package survives unchanged.
\end{enumerate}

So \(\tau\) looks like a good local test object, not a proved extension. The
later scattering attack suggests a cleaner target object here too: not merely a
real channel, but a completed-object package whose background-subtracted zero
part plays the role of \(S(m)\). A later quotient-phase attack sharpens this in
the same narrow way: the self-dual quotient is the cleanest \emph{phase}
candidate, but none of the five localization checkpoints is discharged by that
alone.

\section{The Missing Realization Theorem}

The main missing theorem suggested by the cycle is a completed-\(L\)
realization theorem with hypotheses of the following kind:

\begin{enumerate}
\item a real local phase on the critical line;
\item a background/value split \(q=B+S\);
\item a real, sign-compatible amplitude replacing \(S(m)\);
\item a corrected odd holomorphic transverse scalar on a microscopic disk;
\item denominator comparability and tail control on that disk;
\item preserved whitening scales strong enough to recover the same
\(Q^{-3}\)-type boundary hierarchy.
\end{enumerate}

Only after that theorem exists does it become reasonable to discuss GRH-style
extensions of the current package. The follow-up attack also shows that one
should not describe the calibration layer as source-light from local scope
alone.

\section{What Can Be Ported Back Immediately}

Even before any completed-\(L\) realization theorem is written, the main RH
draft would benefit from a sharper internal boundary:

\begin{enumerate}
\item present the opening local package in neutral phase language;
\item mark the first zeta-side input layer explicitly;
\item present the \(N\)-point projector as universal on odd holomorphic germs;
\item keep the final contradiction and provenance remarks explicitly zeta-side.
\end{enumerate}

This is the strongest safe takeaway from the present cycle.

\section{Near-Term Program In \texttt{grh/}}

The next steps suggested by the cycle are:

\begin{enumerate}
\item a family-agnostic note isolating the local phase package and its minimal
hypotheses;
\item a Dirichlet note that decides the correct critical-line object and tracks
where sign-compatibility of the amplitude is used;
\item a \(\tau\) note that checks the degree-\(2\) archimedean background,
microscopic radius, and whitening power counting.
\end{enumerate}

\end{document}
